Interstellar dust is heated up by hot ionizing radiation from young blue stars. They then re-emit the UV light.
% sic: The official wording says re-emit UV, but a 50 K blackbody peaks in the infrared, consistent with the official 57.96 micrometre answer.
A dust cloud in a dusty galaxy was recently studied by a researcher. Make the simplifying approximation that you can treat the dust cloud as a blackbody. Help him to answer the following questions.
\begin{enumerate}
\item Find the wavelength where the dust emission peaks $\lambda_{\max}$ if it has a temperature of $T=50\,\mathrm{K}$.
\item The researcher determines that the dust has a total luminosity $L_{\mathrm{dust}}=1.6\times10^{11}L_\odot$. Estimate the size of the dust cloud, assuming a spherical distribution.
\item If the cosmological redshift of the galaxy $z$ is 0.01, estimate the angular size $\theta$ of the dust cloud.
\item What aperture telescope is minimally needed if we want to resolve the dust cloud?
\end{enumerate}
\textit{Supplied examination data.} $\lambda_mT=2.898\times10^{-3}\,\mathrm{m\,K}$, $L_\odot=3.826\times10^{26}\,\mathrm{W}$, $T_{\mathrm{eff},\odot}=5778\,\mathrm{K}$, $R_\odot=695500\,\mathrm{km}$, $c=2.998\times10^8\,\mathrm{m\,s^{-1}}$, $H_0=70\,\mathrm{km\,s^{-1}\,Mpc^{-1}}$, and $1\,\mathrm{pc}=3.086\times10^{16}\,\mathrm{m}$.
